\section{章节罗盘：从入关胜利到可治理的帝国}
\label{sec:early-qing-compass}

本章不把1644年写成一个已经完成的统一时刻。山海关、北京与南方战场之间，清廷必须把八旗的军事突破转化为能够收粮、发令、任官和维持秩序的朝廷；而降官、地方士绅、驻防军与居民进入这一秩序的条件并不相同。征服的速度、暴力与行政接收将并列呈现，而不会被一个“定鼎”概念遮蔽。

叙事随后转向摄政、亲政、撤藩、海疆和北部、西部边地。这里的关键问题不是中央是否发布命令，而是兵饷、迁徙、册封、设治和当地中介如何使命令在不同地区取得有限效力。本章结尾将把清初的整合看作一份持续支出的政治账：它提供了盛世帝国可以扩张的框架，也留下了驻军、身份和边疆供给的长期成本。
